\documentclass[12pt,a4paper]{article}

\usepackage[T1]{fontenc}
\usepackage[utf8]{inputenc}
\usepackage{lmodern}
\usepackage[left=30mm,right=25mm,top=25mm,bottom=25mm]{geometry}
\usepackage[english, ngerman]{babel}
\usepackage[autostyle=true,german=quotes]{csquotes}
\usepackage[onehalfspacing]{setspace}
\usepackage[expansion=false]{microtype}
\usepackage{booktabs}
\usepackage{tabularx}
\usepackage{enumitem}
\usepackage{xcolor}
\usepackage{amsmath}
\usepackage{etoolbox}
\usepackage{caption}
\usepackage{tikz}
\usetikzlibrary{positioning,patterns}
\usepackage{pgfplots}
\pgfplotsset{compat=1.18}
% Diagrammfarben (geprüft auf Farbsehschwächen), zusätzlich Marker/Strichart
\definecolor{cblue}{HTML}{2A78D6}
\definecolor{corange}{HTML}{EB6834}
\definecolor{ctext}{HTML}{52514E}
\usepackage[backend=biber,style=ieee,urldate=long]{biblatex}
\addbibresource{sycl.bib}
\usepackage{hyperref}

\hypersetup{
    colorlinks,
    citecolor   = black,
    filecolor   = black,
    linkcolor   = black,
    urlcolor    = black,
    pdftitle    = {Von C++ zu SYCL: Aufwand, Leistung und Portabilität},
    pdfauthor   = {Erlind Sejdiu},
    pdfcreator  = {pdflatex},
    pdfproducer = {LaTeX with hyperref}
}

% Tabellen, Abbildungen und Literaturverzeichnis einzeilig
\AtBeginEnvironment{tabular}{\singlespacing}
\AtBeginEnvironment{tabularx}{\singlespacing}
\AtBeginEnvironment{tikzpicture}{\singlespacing}
\captionsetup{font=small, labelfont=bf}
\AtBeginBibliography{\singlespacing}

% Markierung für Stellen, die noch ergänzt/geprüft werden müssen
\newcommand{\offen}[1]{\textcolor{red}{[#1]}}

\begin{document}

% ---------------------------------------------------------------------------
% Titelei (römische Seitenzahlen, Titelblatt ohne Nummer)
% ---------------------------------------------------------------------------
\pagenumbering{Roman}

\begin{titlepage}
    \raggedright
    Erlind Sejdiu\\
    106357
    \vspace*{3cm}

    \begin{center}
        {\LARGE\bfseries Von C++ zu SYCL\\[0.3em]
        Umstellungsaufwand, Leistung und Portabilität\\[0.3em]
        am Beispiel der Matrixmultiplikation\par}
        \vspace{2cm}
        {\large Seminararbeit\\
        der Hochschule Karlsruhe\\
        im Master Studiengang Informatik\\
        im Wintersemester 2026/27\par}
    \end{center}
    \vfill
    \begin{tabular}{@{}ll}
        Betreuer: & \offen{Name der betreuenden Person}\\
        Themengruppe: & \offen{Themengruppe des Seminars}\\
        Abgabedatum: & \offen{Datum}
    \end{tabular}
\end{titlepage}

\setcounter{page}{1}
\tableofcontents
\newpage
\listoffigures
\listoftables
\newpage

% ---------------------------------------------------------------------------
% Hauptteil (arabische Seitenzahlen ab 1)
% ---------------------------------------------------------------------------
\pagenumbering{arabic}

\section{Einleitung}

Moderne Rechner sind heterogen: Neben der CPU stehen GPUs zur Verfügung, die
bei datenparallelen Aufgaben um Größenordnungen schneller rechnen. Ihre
Programmierung ist jedoch herstellerspezifisch. CUDA läuft nur auf
NVIDIA-Hardware \cite{nickolls2008cuda}, OpenCL ist zwar offen, trennt aber
Host- und Kernel-Code in verschiedene Sprachen \cite{stone2010opencl}, und
Bibliotheken wie Kokkos \cite{edwards2014kokkos} oder RAJA
\cite{beckingsale2019raja} verlagern das Problem in eine zusätzliche
Abstraktionsschicht. SYCL verspricht, mit reinem C++ und einem einzigen
Quellcode CPUs und GPUs verschiedener Hersteller zu nutzen
\cite{khronos2020sycl}. Diese Arbeit prüft das Versprechen anhand der
Leitfrage:

\begin{quote}
\emph{Wie aufwendig ist die Umstellung eines C++-Programms auf SYCL, welche
Leistung bringt sie, und gilt die Portabilität auch für die Leistung?}
\end{quote}

\noindent Dazu wird ein sequentielles C++-Programm zur Matrixmultiplikation in
mehreren Stufen nach SYCL überführt und jede Stufe auf einer CPU und einer
NVIDIA-GPU vermessen.

% ---------------------------------------------------------------------------
\section{Grundlagen von SYCL}\label{sec:grundlagen}
% ---------------------------------------------------------------------------

\paragraph{Standard und Single-Source.}
SYCL ist eine Spezifikation der Khronos Group, kein Compiler. Seit SYCL~2020
baut sie auf C++17 auf und ist nicht mehr an OpenCL gebunden
\cite{khronos2020sycl}. Implementierungen wie DPC++ von Intel
\cite{intel_llvm}, das auch Teil der oneAPI-Werkzeuge ist, und AdaptiveCpp
(früher hipSYCL) der Universität Heidelberg \cite{alpay2020hipsycl} bilden das Programm auf Backends wie CUDA, HIP,
Level~Zero, OpenCL oder OpenMP ab (Abbildung~\ref{fig:stack}a). Grundprinzip
ist \emph{Single-Source}: Der Kernel, also der Gerätecode, ist ein
gewöhnliches C++-Lambda in derselben Datei wie das Hauptprogramm. Er wird
beim Übersetzen in eine Zwischendarstellung gebracht und beim ersten Start
\emph{just in time} für das konkrete Gerät kompiliert
\cite{alpay2023sscp}. Im Kernel sind Ausnahmen, virtuelle Aufrufe, Rekursion
und dynamischer Speicher verboten, und Daten werden per Wertkopie übernommen.
Weil Host- und Gerätecode in einer Datei stehen, kann der Compiler Typen,
Templates und Hilfsfunktionen auf beiden Seiten gemeinsam prüfen. Anders als
bei OpenCL entfallen damit getrennte Kernel-Dateien und das Übergeben von
Kernel-Argumenten über eine eigene Schnittstelle, anders als bei CUDA werden
keine Spracherweiterungen benötigt.

\begin{figure}[htbp]
\centering
\begin{minipage}[b]{0.47\textwidth}\centering
\begin{tikzpicture}[
    box/.style={draw, minimum width=6.4cm, minimum height=0.55cm, align=center, fill=gray!10, font=\scriptsize},
    node distance=0.3cm
]
    \node[box] (src) {C++17-Quellcode mit SYCL};
    \node[box, below=of src] (impl) {Implementierung: DPC++ \,|\, AdaptiveCpp};
    \node[box, below=of impl] (be) {Backend: OpenCL \,|\, CUDA \,|\, HIP \,|\, Level Zero \,|\, OpenMP};
    \node[box, below=of be] (hw) {Hardware: CPU \,|\, NVIDIA- \,|\, AMD- \,|\, Intel-GPU};
    \draw[->] (src) -- (impl);
    \draw[->] (impl) -- (be);
    \draw[->] (be) -- (hw);
\end{tikzpicture}\\[0.3em]
{\small (a) Schichtenmodell}
\end{minipage}\hfill
\begin{minipage}[b]{0.5\textwidth}\centering
\begin{tikzpicture}[scale=0.7, font=\scriptsize]
    \fill[cblue!15] (1,2) rectangle (2,3);
    \draw[step=0.25, gray!60] (1,2) grid (2,3);
    \fill[cblue] (1.25,2.5) rectangle (1.5,2.75);
    \draw[step=1, gray] (0,0) grid (4,4);
    \draw[thick] (0,0) rectangle (4,4);
    \draw[thick, cblue] (1,2) rectangle (2,3);
    \draw[->] (4.3,3.4) node[right, align=left] {Work-Group\\($T \times T$)} -- (2.05,2.85);
    \draw[->] (4.3,1.6) node[right, align=left] {Work-Item} -- (1.45,2.55);
    \node[below] at (2,0) {Indexraum der \texttt{nd\_range} ($n \times n$)};
\end{tikzpicture}\\[0.3em]
{\small (b) Gliederung des Indexraums}
\end{minipage}
\caption{Aufbau von SYCL und Ausführungsmodell}
\label{fig:stack}
\end{figure}

\paragraph{Queues und Geräte.}
Der Host, meist die CPU, reicht Arbeit über eine \emph{Queue} an genau ein
\emph{Gerät} ein. Welches Gerät das ist, bestimmt ein Device Selector oder
von außen eine Umgebungsvariable, ohne den Quellcode zu ändern. Über die
Abfrageschnittstelle \texttt{get\_info} kann ein Programm zur Laufzeit
Eigenschaften des Geräts wie Name, Anzahl der Recheneinheiten oder maximale
Work-Group-Größe auslesen und sich daran anpassen. Die
Ausführung ist asynchron: Jeder Auftrag liefert ein \emph{Event}, auf das der
Host warten kann. Die Laufzeit baut aus den Aufträgen einen
Abhängigkeitsgraphen und kann voneinander unabhängige Transfers und
Berechnungen überlappen, ohne dass der Programmierer Threads verwaltet
\cite{reinders2023dpcpp}.

\paragraph{Ausführungsmodell.}
Ein Kernel wird mit \texttt{parallel\_for} für jeden Punkt eines Indexraums
einmal ausgeführt; eine Ausführung heißt \emph{Work-Item}. Bei einer
\texttt{nd\_range} werden Work-Items zu \emph{Work-Groups} gebündelt, die
gemeinsam auf einer Recheneinheit laufen, schnellen lokalen Speicher teilen
und sich über Barrieren synchronisieren (Abbildung~\ref{fig:stack}b). Auf NVIDIA-GPUs entspricht eine
Work-Group einem Thread-Block mit höchstens 1024 Threads
\cite{nvidia_cuda_guide}, auf der CPU wird sie von einem Kern vektorisiert
abgearbeitet. Seit SYCL~2020 gibt es zusätzlich \emph{Sub-Groups}, die auf der
GPU einem Warp und auf der CPU den SIMD-Lanes entsprechen
\cite{khronos2020sycl}.

\paragraph{Speichermodell.}
Daten zwischen Host und Gerät werden entweder über Buffer und Accessoren oder
über zeigerbasiertes \emph{Unified Shared Memory} (USM) verwaltet. Bei Buffern
leitet die Laufzeit aus den deklarierten Zugriffen ab, wann sie Daten kopieren
muss; USM entspricht dagegen dem gewohnten Umgang mit Zeigern in C++. Bei USM
\emph{shared} verschiebt die Laufzeit die Daten automatisch, bei USM
\emph{device} kopiert der Programmierer selbst. Auf dem Gerät unterscheidet
SYCL zwischen globalem Speicher, der für alle Work-Items sichtbar, groß, aber
langsam ist, lokalem Speicher, den sich die Work-Items einer Work-Group teilen
und der klein, aber schnell ist, sowie privatem Speicher in den Registern
eines Work-Items \cite{reinders2023dpcpp}. Diese Hierarchie gezielt zu
nutzen ist, wie bei CUDA, der wichtigste Hebel für Leistung
\cite{hwu2022pmpp}.

% ---------------------------------------------------------------------------
\section{Vorgehen}\label{sec:vorgehen}
% ---------------------------------------------------------------------------

\subsection{Beispielproblem}

Untersucht wird die Multiplikation $C = A \cdot B$ zweier $n \times
n$-Matrizen in einfacher Genauigkeit. Jeder Eintrag $C_{ij}$ ist das
Skalarprodukt der Zeile $i$ von $A$ mit der Spalte $j$ von $B$. Die Rechnung
benötigt $2n^3$ Gleitkommaoperationen bei nur $3n^2$ Werten Datenvolumen, und
alle Einträge von $C$ sind voneinander unabhängig. Damit ist das Problem gut
parallelisierbar, rechenintensiv und leicht auf Korrektheit prüfbar. Es gilt
als Standardbeispiel der GPU-Programmierung \cite{hwu2022pmpp} und ist seit
Langem Gegenstand der Leistungsoptimierung auf GPUs \cite{volkov2008gemm}.

\subsection{Stufenmodell}

Ein sequentielles C++-Programm wird schrittweise nach SYCL überführt. Jede
Stufe führt genau ein neues Konzept ein (Tabelle~\ref{tab:stufen}). So lässt
sich der Umstellungsaufwand pro Konzept beziffern und die Wirkung jedes
Schritts getrennt messen. Die Reihenfolge folgt dem Weg, den ein Entwickler
bei der Portierung typischerweise geht: zuerst die Parallelisierung mit
möglichst wenig Änderung, dann die Kontrolle über den Speicher und zuletzt
die hardwarenahe Optimierung. Stufe~4 enthält bewusst keinen neuen Code: Sie
prüft, ob die Programme der Stufen~1 bis~3 unverändert auf einem anderen
Gerät laufen. Alle Stufen verwenden dieselben Hilfsfunktionen für
Initialisierung, Zeitmessung und Ergebnisprüfung, sodass sich ihre
Messwerte direkt vergleichen lassen.

\begin{table}[htbp]
\centering
\small
\begin{tabularx}{\textwidth}{@{}c>{\raggedright\arraybackslash}X>{\raggedright\arraybackslash}X@{}}
\toprule
\textbf{Stufe} & \textbf{Änderung} & \textbf{Neue Konzepte} \\
\midrule
0 & Ausgangsprogramm, drei Schleifen, ein CPU-Kern & -- \\
1 & Äußere Schleifen durch \texttt{parallel\_for} ersetzt & Queue, Kernel, \texttt{range}, USM \emph{shared} \\
2 & Getrennter Gerätespeicher, explizites Kopieren & USM \emph{device}, \texttt{memcpy} \\
3 & Kachelung im lokalen Speicher & \texttt{nd\_range}, Work-Group, Barriere \\
4 & Kein neuer Code, Stufen 1--3 auf anderem Gerät & Geräteauswahl, Backends \\
\bottomrule
\end{tabularx}
\caption{Stufen der Umstellung}
\label{tab:stufen}
\end{table}

\subsection{Bewertung der Performance-Portabilität}\label{sec:pp}

Für die Bewertung werden drei Ebenen getrennt: \emph{funktionale
Portabilität} (das Programm läuft überall korrekt),
\emph{Implementierungseffizienz} (wie gut eine Variante ein einzelnes Gerät
ausnutzt) und \emph{Performance-Portabilität} (ob ein ähnlich hoher Anteil der
erreichbaren Leistung auf allen Geräten erzielt wird). Pennycook, Sewall und Lee messen dies als harmonischen
Mittelwert der Effizienzen $e_i$ über eine Plattformmenge $H$
\cite{pennycook2016metric}:
\[
    \mathrm{PP}(a, p, H) =
    \begin{cases}
        |H| \,\big/ \sum_{i \in H} 1/e_i(a,p) & \text{falls $a$ auf allen $i \in H$ läuft,}\\
        0 & \text{sonst.}
    \end{cases}
\]
Dabei ist $a$ die Anwendung und $p$ das Problem. Als Effizienz wird hier die
\emph{Anwendungseffizienz} verwendet, also die Leistung relativ zur besten
gemessenen Variante auf dem jeweiligen Gerät. Eine einzelne schwache
Plattform senkt den Wert stark, und der Wert gilt nur für die gewählte
Plattformmenge \cite{pennycook2019implications}. Frühere Studien fanden
SYCL auf vielen Plattformen konkurrenzfähig, aber deutlich abhängig von
Implementierung und Parametrierung \cite{deakin2020sycl, breyer2022svm}.
Diese Arbeit ergänzt sie um eine kleine Fallstudie, die nicht fertige
Implementierungen vergleicht, sondern misst, welcher Schritt der Umstellung
welchen Effekt auf welchem Gerät hat.

\subsection{Messumgebung und Methodik}

Gemessen wird auf einem AMD Ryzen~5 5600X mit sechs Kernen und einer NVIDIA
GeForce RTX~5060~Ti. Als SYCL-Compiler dient DPC++ in Version~7.2 aus dem
Open-Source-Repository \cite{intel_llvm}. Die CPU wird über das
OpenCL-Backend, die GPU über das CUDA-Backend angesprochen. Stufe~0 wird mit
\texttt{g++ -O3} übersetzt.

Jede Stufe wird für $n \in \{256, 512, 1024, 2048\}$ ausgeführt. Der erste
Durchlauf wird getrennt erfasst, weil er JIT-Kompilierung und
Geräteinitialisierung enthält. Danach folgen fünf Messläufe, von denen der
Median berichtet wird. Bei den Stufen~2 und~3 werden Kernel- und Transferzeit
getrennt gemessen. Jedes Ergebnis wird an 64 zufälligen Einträgen gegen eine
Referenzrechnung in doppelter Genauigkeit geprüft; die zulässige Abweichung
wächst mit $\sqrt{n}$, um Rundungsfehler in einfacher Genauigkeit zu
berücksichtigen. Die Leistung wird in GFLOPS angegeben, also $2n^3$ geteilt
durch die Kernel-Zeit.

Die Matrizen werden mit festen Startwerten zufällig im Intervall $[-1, 1]$
belegt, sodass jeder Lauf dieselben Eingaben verarbeitet. Gemessen wird mit
\texttt{std::chrono::steady\_clock} auf dem Host, wobei jede Messung erst
nach dem Warten auf das Ende des Kernels endet. Der Median ist gegenüber
einzelnen Ausreißern, etwa durch Hintergrundprozesse, robuster als der
Mittelwert. Sämtliche Messungen werden über Skripte automatisiert und als
CSV-Dateien abgelegt.

% ---------------------------------------------------------------------------
\section{Umsetzung}\label{sec:umsetzung}
% ---------------------------------------------------------------------------

\subsection{Werkzeugkette}

Die Standardinstallation von Intel oneAPI erkannte auf dem Testsystem nur die
CPU. Für die NVIDIA-GPU war ein eigener Build von DPC++ mit CUDA-Unterstützung
nötig. Beim Übersetzen müssen alle Zielarchitekturen angegeben werden
(\texttt{-fsycl-targets}); ein nur für NVIDIA übersetztes Programm bricht auf
der CPU zur Laufzeit ab, weil kein passendes Geräteabbild vorhanden ist. Der
Einrichtungsaufwand war damit größer als der Programmieraufwand der ersten
Stufe. Für die Stufen~1 bis~3 werden die Programme deshalb einmal für beide
Ziele zugleich übersetzt, sodass dieselbe ausführbare Datei auf CPU und GPU
läuft.

\subsection{Stufe 1: Der erste Kernel}

Gegenüber dem sequentiellen Programm ändern sich drei Dinge: Es wird eine
Queue angelegt, die Matrizen werden als USM \emph{shared} reserviert, und die
beiden äußeren Schleifen über Zeilen und Spalten entfallen zugunsten eines
zweidimensionalen \texttt{parallel\_for} über $n \times n$ Work-Items. Die
innere Schleife und damit die Rechenvorschrift bleiben unverändert; jedes
Work-Item berechnet einen Eintrag $C_{ij}$. Insgesamt betrifft der Umbau etwa
fünf Zeilen. Einzige Stolperstelle ist die Wertübernahme in den Kernel: Die
Matrizen müssen als Zeiger, nicht als Referenzen auf Host-Objekte übergeben
werden.

Die Queue wird ohne Angabe eines Geräts angelegt, die Laufzeit wählt also das
als am leistungsfähigsten eingeschätzte Gerät. Da der Kernel asynchron
startet, wartet der Host anschließend ausdrücklich auf sein Ende. Weil USM
\emph{shared} auf Host und Gerät gültig ist, kann das Ergebnis danach ohne
weiteren Schritt auf dem Host geprüft werden. Der Programmierer muss sich in
dieser Stufe also weder um Datentransfers noch um die Aufteilung der Arbeit
auf die Hardware kümmern.

\subsection{Stufe 2: Expliziter Gerätespeicher}

Die Matrizen werden als USM \emph{device} angelegt. Die Eingaben werden vor
dem Kernel ausdrücklich auf das Gerät und das Ergebnis danach zurück kopiert.
Der Kernel bleibt identisch. Ziel dieser Stufe ist Transparenz: Die
Transferkosten, die bei USM \emph{shared} unsichtbar anfallen, werden
getrennt messbar. Zugleich gewinnt der Programmierer Kontrolle darüber, wann
Daten bewegt werden. In einer realen Anwendung ließen sich so Eingaben
einmalig kopieren und für mehrere Kernel wiederverwenden oder Transfers über
die asynchrone Queue mit Berechnungen überlappen.

\subsection{Stufe 3: Kachelung im lokalen Speicher}

In den Stufen~1 und~2 liest jedes Work-Item eine ganze Zeile von $A$ und eine
ganze Spalte von $B$ aus dem globalen Speicher. Benachbarte Work-Items lesen
dabei immer wieder dieselben Werte. Die etablierte Gegenmaßnahme ist die
Kachelung \cite{hwu2022pmpp, volkov2008gemm}. Der Kernel wird dazu mit einer
\texttt{nd\_range} gestartet, sodass Work-Groups von $T \times T$ Work-Items
entstehen. Wie Abbildung~\ref{fig:tiling} zeigt, lädt jede Work-Group in einer
Schleife gemeinsam je einen $T \times T$-Block von $A$ und $B$ in den lokalen
Speicher, wartet an einer Barriere, bis alle Werte geladen sind, rechnet damit
weiter und wartet an einer zweiten Barriere, bevor die nächsten Blöcke die
alten überschreiben.

\begin{figure}[htbp]
\centering
\begin{tikzpicture}[font=\scriptsize]
    \def\s{2.4}
    % A
    \fill[cblue!15] (0,1.2) rectangle (\s,1.6);
    \fill[cblue] (0.4,1.2) rectangle (0.8,1.6);
    \draw[step=0.4, gray!60] (0,1.2) grid (\s,1.6);
    \draw (0,0) rectangle (\s,\s);
    \draw[->, thick] (0.6,1.0) -- (1.8,1.0);
    \node[above] at (\s/2,\s) {$A$};
    \node[below, align=center] at (\s/2,-0.05) {Zeilenband,\\Block wandert $\rightarrow$};
    \node at (\s+0.45,\s/2) {\large$\cdot$};
    % B
    \begin{scope}[xshift=3.3cm]
    \fill[cblue!15] (1.2,0) rectangle (1.6,\s);
    \fill[cblue] (1.2,1.6) rectangle (1.6,2.0);
    \draw[step=0.4, gray!60] (1.2,0) grid (1.6,\s);
    \draw (0,0) rectangle (\s,\s);
    \draw[->, thick] (1.8,1.8) -- (1.8,0.6);
    \node[above] at (\s/2,\s) {$B$};
    \node[below, align=center] at (\s/2,-0.05) {Spaltenband,\\Block wandert $\downarrow$};
    \node at (\s+0.45,\s/2) {\large$=$};
    \end{scope}
    % C
    \begin{scope}[xshift=6.6cm]
    \fill[corange] (1.2,1.2) rectangle (1.6,1.6);
    \draw (0,0) rectangle (\s,\s);
    \node[above] at (\s/2,\s) {$C$};
    \node[below, align=center] at (\s/2,-0.05) {Ergebnisblock\\einer Work-Group};
    \end{scope}
\end{tikzpicture}
\caption{Kachelung: Die dunklen Blöcke liegen gerade im lokalen Speicher}
\label{fig:tiling}
\end{figure}

Jeder Wert wird so pro Work-Group nur einmal statt $T$-mal aus dem globalen
Speicher gelesen. Die Stufe setzt deutlich mehr Hardwareverständnis voraus.
Fehlt etwa die zweite Barriere, überschreiben schnelle Work-Items den Block,
während langsamere noch daraus lesen. Das führt nicht zum Absturz, sondern zu
sporadisch falschen Ergebnissen.

Zwei Details sind zu beachten. Erstens muss der globale Indexraum einer
\texttt{nd\_range} ein Vielfaches der Work-Group-Größe sein. Er wird daher
auf das nächste Vielfache von $T$ aufgerundet, und Werte außerhalb der Matrix
werden beim Laden durch null ersetzt. Zweitens belegt jede Work-Group
$2 \cdot T^2$ Gleitkommazahlen im lokalen Speicher, bei $T = 16$ also 2\,KiB
und bei $T = 64$ bereits 32\,KiB. Die Kachelgröße wirkt sich damit nicht nur
auf die Wiederverwendung der Daten aus, sondern auch darauf, wie viele
Work-Groups gleichzeitig auf einer Recheneinheit Platz finden.

\subsection{Stufe 4: Geräteauswahl und Kachelgrößen}

Für Stufe~4 werden dieselben Programme nur mit anderer Umgebungsvariable
gestartet, \texttt{ONEAPI\_DEVICE\_SELECTOR=opencl:cpu} für die CPU und
\texttt{cuda:gpu} für die GPU. Um die Performance-Portabilität zu bewerten,
wird Stufe~3 zusätzlich mit den Kachelgrößen $T \in \{8, 16, 32, 64\}$
übersetzt. Die Kachelgröße ist ein Übersetzungsparameter, der Quellcode
bleibt also gleich. Als Standardwert dient $T = 16$, weil damit eine
Work-Group 256 Work-Items umfasst, eine auf GPUs übliche Größe. Jede Variante wird auf beiden Geräten für $n = 1024$ und
$n = 2048$ gemessen. Die Effizienzen und die PP-Kennzahl werden
anschließend per Skript aus den Messdaten berechnet.

% ---------------------------------------------------------------------------
\section{Ergebnisse und Diskussion}\label{sec:ergebnisse}
% ---------------------------------------------------------------------------

\subsection{Leistung, Anlauf und Transfer}

Alle Stufen rechneten auf beiden Geräten korrekt. Tabelle~\ref{tab:leistung}
zeigt die Kernel-Zeiten für $n = 1024$. Schon Stufe~1 erreicht auf der GPU
rund 80\,\% der Leistung von Stufe~3. Der Faktor von etwa 3000 gegenüber
Stufe~0 ist kein fairer CPU-GPU-Vergleich, da Stufe~0 nur einen Kern nutzt.
Die Kachelung wirkt je Gerät sehr unterschiedlich: auf der GPU rund 30\,\%,
auf der CPU ein Faktor~4,5. Vermutlich fängt die GPU redundante Zugriffe
schon über Caches und gebündelte Speicherzugriffe ab, während die CPU erst
durch die Kachelung cachefreundlich wird \cite{hwu2022pmpp}. Bei
$n = 256$ ist Stufe~3 auf der GPU sogar etwas langsamer als Stufe~2, weil der
Aufwand für die Barrieren bei so wenig Arbeit nicht aufgewogen wird. 

\begin{table}[htbp]
\centering
\small
\begin{tabular}{@{}lrr@{}}
\toprule
\textbf{Stufe ($n = 1024$)} & \textbf{CPU} & \textbf{GPU} \\
\midrule
0 sequentiell (1 Kern) & 4484\,ms & -- \\
1 USM \emph{shared}    & 87{,}0\,ms & 1{,}94\,ms (1105 GFLOPS) \\
2 USM \emph{device}    & 83{,}1\,ms & 1{,}97\,ms \\
3 Kachelung, $T = 16$  & 18{,}4\,ms & 1{,}54\,ms (1395 GFLOPS) \\
\bottomrule
\end{tabular}
\caption{Kernel-Zeit ohne Transfer, Median aus fünf Läufen}
\label{tab:leistung}
\end{table}

Der erste Lauf ist durch die JIT-Kompilierung deutlich langsamer: bei
$n = 256$ auf der GPU 0,75\,ms statt 0,04\,ms, auf der CPU sogar 118\,ms statt
0,4\,ms. Stufe~1 und~2 sind im Median
gleich schnell (1,94 gegenüber 1,97\,ms), da USM \emph{shared} die Daten nur
beim ersten Zugriff verschiebt.  Bei kleinen Matrizen dominiert
zudem der Transfer (Abbildung~\ref{fig:ergebnisse}a), der bei $n = 256$ 82\,\% und bei $n = 2048$ nur noch
20\,\% der Gesamtzeit ausmacht, weil die Rechnung mit $n^3$, die Datenmenge
aber nur mit $n^2$ wächst. Ob sich die GPU lohnt, hängt also von der
Problemgröße ab.

\subsection{Performance-Portabilität}

Abbildung~\ref{fig:ergebnisse}b zeigt die Effizienz je Kachelgröße für
$n = 2048$ relativ zum besten Wert des Geräts. Das Optimum ist
geräteabhängig: $T = 16$ auf der GPU, $T = 32$ auf der CPU. Bei $T = 64$
umfasst eine Work-Group 4096 Work-Items, mehr als NVIDIA erlaubt; der Kernel
startet dort nicht, und die PP ist null, obwohl die CPU 91\,\% erreicht. Mit
$T = 32$ gibt es für diese zwei Geräte einen guten Kompromiss
(PP~0,99, bei $n = 1024$ 0,96), der sich aber nur durch Messen auf allen
Geräten finden lässt. Eine falsche Wahl trifft die Geräte zudem ungleich:
$T = 8$ kostet die GPU 17\,\%, die CPU 55\,\% ihrer Leistung.

\begin{figure}[htbp]
\centering
\pgfplotsset{
    every axis/.append style={
        font=\scriptsize, width=7cm, height=4.8cm,
        axis line style={ctext}, tick style={ctext},
        grid style={gray!20}, ymajorgrids,
        legend style={draw=none, fill=none, font=\scriptsize},
        legend cell align=left,
    }
}
\begin{minipage}[t]{0.49\textwidth}\centering
\begin{tikzpicture}
\begin{axis}[
    xmode=log, ymode=log, log basis x=2,
    xtick={256,512,1024,2048}, xticklabels={256,512,1024,2048},
    xlabel={Matrixgröße $n$}, ylabel={Zeit in ms},
    legend pos=north west,
    title={(a) Rechnung und Transfer, Stufe 2, GPU},
]
\addplot[cblue, line width=1.2pt, mark=*, mark size=2pt]
    coordinates {(256,0.038) (512,0.263) (1024,1.974) (2048,15.869)};
\addlegendentry{Rechnung}
\addplot[corange, line width=1.2pt, dashed, mark=square*, mark size=2pt, mark options={solid}]
    coordinates {(256,0.173) (512,0.405) (1024,1.344) (2048,3.934)};
\addlegendentry{Transfer}
\end{axis}
\end{tikzpicture}
\end{minipage}\hfill
\begin{minipage}[t]{0.49\textwidth}\centering
\begin{tikzpicture}
\begin{axis}[
    ybar, bar width=7pt, enlarge x limits=0.18,
    ymin=0, ymax=115, ytick={0,25,50,75,100},
    symbolic x coords={8,16,32,64}, xtick=data,
    xticklabels={{$T{=}8$\\PP 0{,}59},{$T{=}16$\\PP 0{,}93},{$T{=}32$\\PP 0{,}99},{$T{=}64$\\PP 0}},
    xticklabel style={align=center},
    ylabel={Effizienz in \%},
    legend style={at={(0.5,1.0)}, anchor=south, legend columns=2},
    title style={yshift=3ex},
    title={(b) Effizienz je Kachelgröße, $n = 2048$},
]
\addplot[fill=cblue, draw=white, line width=0.8pt]
    coordinates {(8,83) (16,100) (32,99) (64,0)};
\addlegendentry{GPU}
\addplot[fill=corange, draw=white, line width=0.8pt, postaction={pattern=north east lines, pattern color=white}]
    coordinates {(8,45) (16,88) (32,100) (64,91)};
\addlegendentry{CPU}
\node[font=\scriptsize, text=ctext, rotate=90, anchor=west] at ([xshift=-4.5pt]axis cs:64,2) {startet nicht};
\end{axis}
\end{tikzpicture}
\end{minipage}
\caption{Transferanteil (a) und Performance-Portabilität der Kachelgrößen (b)}
\label{fig:ergebnisse}
\end{figure}

\subsection{Diskussion und Fazit}

Der Einstieg in SYCL ist klein: Wenige Zeilen genügen für einen Großteil der
GPU-Leistung, und derselbe Code läuft ohne Änderung auf CPU und GPU. Die
funktionale Portabilität ist damit erreicht. Aufwendig wurden erst die
Werkzeugkette und hardwarenahe Optimierungen, deren Wissen aber auch für CUDA
nötig ist. Performance-Portabilität stellt sich dagegen nicht automatisch
ein: SYCL abstrahiert die Schnittstelle, nicht die Hardware. Parameter wie
die Kachelgröße müssen pro Gerät bestimmt werden, und eine ungünstige Wahl
kann ein Gerät ganz ausschließen. Das bestätigt im Kleinen die Ergebnisse von
\cite{deakin2020sycl, breyer2022svm}.

Die Beweiskraft der Fallstudie ist allerdings begrenzt. Erstens fehlt eine
starke \emph{Baseline}: Stufe~0 ist bewusst naiv, und die Effizienzen
beziehen sich auf die beste eigene Variante statt auf optimierte Bibliotheken
wie cuBLAS oder oneMKL. Zweitens gilt der PP-Wert
nur für die gewählte Menge aus zwei Geräten \cite{pennycook2019implications}.
Drittens ist die dichte Matrixmultiplikation rechenintensiv und regelmäßig,
was nicht für alle Anwendungen gilt. Viertens wurden nur ein Compiler und eine Messreihe
untersucht, sodass Effekte der Toolchain nicht von denen der Hardware
getrennt sind und kleine Unterschiede im Messrauschen liegen können
\cite{georges2007rigorous}. Ein belastbareres Versuchsdesign würde daher
weitere Geräte und Hersteller, mindestens zwei SYCL-Implementierungen,
Bibliotheken als Referenz, Kernel mit unterschiedlichem Rechen-zu-Speicher-%
Verhältnis und wiederholte Messreihen mit Konfidenzintervallen umfassen.

Im untersuchten Rahmen hält das Versprechen \enquote{einmal schreiben,
überall ausführen} für die Korrektheit, nicht aber für die Leistung.

\newpage
\printbibliography[heading=bibintoc, title={Literaturverzeichnis}]

\end{document}
